\section{Discussion and Limitations}\label{sec:discussion}
The architecture makes three separations explicit: mass change from ingestion, physical device ownership from query identity, and language generation from backend facts. These separations expose errors that could otherwise disappear inside a daily total. However, preserving an error's provenance does not prevent the error. Confirmation can be ignored, a reference portion can be inappropriate, and a correct tool result can be misrepresented in the final response.

The fixed setting limits coverage: intake outside the scales requires manual entry. Cup mass cancellation assumes a consistent vessel, and changes in mounting, vibration or temperature may degrade calibration. Multi-user switching introduces a coordination step that must be included in the usability comparison. A three-person pilot cannot establish general adoption or clinical benefit. Nutrition estimates are constrained by food identification, preparation state, source coverage and portion assumptions.

Tool grounding reduces reliance on model arithmetic but leaves tool selection, argument interpretation, timezone handling and explanation quality as separate risks. Prompt injection or malicious food descriptions must not widen access. Security claims require adversarial tests of actual endpoints, credentials and database policies; architecture diagrams alone are insufficient evidence. Goals should be user-configured and not presented as medical prescriptions.

This manuscript reports a design and prospective methodology. The authorized project documentation describes planned integration; no running prototype, deployed interface or formal E1--E5 result is independently verified here. Image-based food estimation, external health-record imports, wellness scores, LINE support and dashboards remain outside the stated MVP evidence. Those extensions should not be used to imply current completeness.

The switch barrier deliberately sacrifices availability during an unacknowledged transition; E3 reports this coverage cost alongside wrong attribution. Three users and two acquisition sessions per scale offer only a feasibility sample. Neither repeated agent calls nor repeated task placements increase the number of independent people. The additional lifecycle, snapshot and scoring rules are proposed contracts: their presence in this manuscript is not evidence that concurrent software, hardware behavior or access policies have passed tests.
